\SourceLeaf{192}
\Needspace{5\baselineskip}
\pdfbookmark[2]{Werth der noch zu zahlenden Beiträge S.}{witwen-beitraege}
\begin{center}\textit{Werth der noch zu zahlenden Beiträge \(S\).}\end{center}
\noindent\hspace*{1.2em} [Regulativ vom 11. October 1833 und 24. November 1846. Jeder Theilnehmer an der Witwenkasse hat postnumerando jährlich am 17. September den Beitrag an den Rechnungsführer zu entrichten, und werden diese Beiträge von Michaelis zu Michaelis gerechnet. — Wenn ein Mitglied der Witwenkasse in der ersten Hälfte des Beitragsjahres, mithin in den Monaten vom October bis incl. März stirbt, so haben die Erben für das betreffende Jahr den Beitrag nicht mehr einzuzahlen; stirbt dagegen ein Mitglied in der zweiten Hälfte des Jahres, so muss von den Erben am 17. September des Sterbejahres noch der volle Beitrag entrichtet werden. — Die Aufkündigung von Seiten der Theilnehmer muss mittelst schriftlicher Erklärung vor dem 17. September des von Michaelis zu Michaelis laufenden Beitrittsjahres geschehen, der an diesem Tage fällig werdende jährliche Beitrag jedoch noch einmal zu voll bezahlt werden; wer diese Frist nicht einhält, muss für das ganze folgende Beitrittsjahr noch Zahlung leisten, bleibt dann aber bis dahin auch noch Mitglied.

Für den Zeitwerth \(S\) der jährlich als Beitrag zu zahlenden Münzeinheit erhält man daher:]
\[
S\mathord{.}fm\mathord{.}Fm=\rho\mathord{.}fm\mathord{.}F(M+\tfrac{1}{2})+\rho\rho\mathord{.}f(m+1)\mathord{.}F(M+\tfrac{3}{2})+.\ .
\]
\noindent Bezeichnet man also
\[
\psi(m,M)\mathord{.}fm\mathord{.}Fm=\rho f(m+1)F(m+1)+\rho\rho f(m+2)F(m+2)+.\ .\ .
\]
\noindent so ist
\[
\begin{aligned}
Q&=\frac{f(m-\tfrac{1}{24})}{fm}\mathord{.}\frac{F(M-\tfrac{3}{8})}{FM}\mathord{.}\rho^{\frac{1}{4}}\psi(m-\tfrac{1}{24},M-\tfrac{3}{8})\\[5pt]
S&=\frac{f(m-1)}{fm}\mathord{.}\frac{F(M-\tfrac{1}{2})}{FM}\mathord{.}\psi(m-1,M-\tfrac{1}{2})
\end{aligned}
\]
\noindent Schreibt man noch
\[
\begin{aligned}
\frac{f(m-1)}{fm}&=gm\\[5pt]
\frac{F(M-1)}{FM}&=GM
\end{aligned}
\]
\noindent so wird
\[
\begin{aligned}
R&=\tfrac{1}{24}\rho^{\frac{5}{4}}+\rho^{\frac{7}{24}}\varphi m\\
Q&=g(m+\tfrac{23}{48})^{\frac{1}{24}}\mathord{.}G(M+\tfrac{5}{16})^{\frac{3}{8}}\mathord{.}\rho^{\frac{1}{4}}\mathord{.}\psi(m-\tfrac{1}{24},M-\tfrac{3}{8})\\
S&=gm\mathord{.}G(M+\tfrac{1}{4})^{\frac{1}{2}}\mathord{.}\psi(m-1,M-\tfrac{1}{2})
\end{aligned}
\]
